% Shared module: conditioning of the split, Friedrichs angle, regularity.
\subsection{How well separated are the groups? The Friedrichs angle}\label{sec:cond}
\Cref{thm:factor} is algebraic: it says that the split exists, not that it is stable. Stability is a
geometric property of the component spaces $M_j=\ker_{\Hs}Q_j$, the solutions of each factor in the
Sobolev space $\Hs(\Omega)$. Let
\[
  \sumop:\ \textstyle\bigoplus_{j=1}^N M_j\to\Hs(\Omega),\qquad \sumop(u_1,\dots,u_N)=\sum_ju_j,
\]
be the \emph{sum operator}, with $\norm{(u_1,\dots,u_N)}^2=\sum_j\norm{u_j}_{\Hs}^2$, and define the
\emph{conditioning constant}
\begin{equation}\label{eq:gammaN}
  \gN=\inf_{0\ne v\perp\ker\sumop}\frac{\norm{\sumop v}^2}{\norm{v}^2}
  =\inf\sigma\big(\sumop^*\sumop\big|_{(\ker\sumop)^\perp}\big).
\end{equation}
It is the worst case, over all tuples of pieces, of the size of the sum relative to the size of the
pieces. Tuples in $\ker\sumop$ (the overlap of \cref{thm:factor}) are excluded. By the triangle and Cauchy--Schwarz inequalities,
$\norm{\sum_ju_j}^2\le\big(\sum_j\norm{u_j}\big)^2\le N\sum_j\norm{u_j}^2$, hence
\[
  0\le\gN\le N .
\]
If the spaces are mutually orthogonal, $\norm{\sum_ju_j}^2=\sum_j\norm{u_j}^2$ and $\gN=1$. At the other end, if $M_2$
approaches $M_1$ one can take $u_2\approx-u_1$: the pieces stay of unit size while their sum becomes arbitrarily
small, so $\gN\to0$. A small $\gN$ is the signature of groups that can almost cancel each other.

For two groups, $\gamma_2=1-\cF$, where the \emph{Friedrichs cosine}
\[
  \cF=\sup\big\{\abs{\ip ab}:\ a\in M_1\ominus(M_1\cap M_2),\ b\in M_2\ominus(M_1\cap M_2),\
  \norm a=\norm b=1\big\}=\cos\thetaF
\]
is the cosine of the smallest nonzero angle $\thetaF$ between the two spaces
(\cref{fig:friedrichs})~\cite{Deu01}. In regression terms, $M_1$ and $M_2$ are two blocks of regressors and $\cF$ is the
largest correlation between a combination of the first block and one of the second: the closer the blocks are to
collinear, the less the data can attribute to each.

\begin{theorem}[Stability of the split]\label{thm:hm}
Every $u\in\ker_{\Hs}P$ has components $u_j\in\Hs$ if and only if $\sum_jM_j$ is closed and dense in
$\ker_{\Hs}P$. Closedness is equivalent to $\gN>0$, and then
\begin{equation}\label{eq:stability}
  \Big(\sum_j\norm{u_j}_{\Hs}^2\Big)^{1/2}\le\gN^{-1/2}\,\norm{u}_{\Hs}
\end{equation}
for the minimum-norm components. If the component spaces are finite-dimensional, as for the sampled
kernels used by the algorithm, the sum is always closed and density holds whenever the sampled kernel of
$P$ is spanned by those of its factors, so stability reduces to $\gN>0$.
\end{theorem}

Density in $\Hs$ is a genuine condition on infinite-dimensional spaces. On convex domains solutions of
$P(\partial)u=0$ are limits of exponential solutions, uniformly on compact subsets~\cite{Mal56,Hor90},
and every exponential solution lies in $\sum_jM_j$; whether this extends to the $\Hs$ norm up to the boundary
for general $P$ we do not claim.

Two warnings come with the theorem. First, $\gN>0$ can fail on infinite-dimensional component spaces: the pieces of a smooth $u$ may then be
one derivative rougher than $u$ itself. Second, for $N\ge3$ groups the pairwise cosines do not
determine $\gN$: two pairs can be well separated while the three spaces are nearly dependent, so the
full Gram operator of $\sumop$ must be used.

\begin{figure}[tbp]
\centering
\begin{tikzpicture}[scale=1.15,>=Stealth]
  \draw[->,gray] (-0.2,0)--(3.8,0); \draw[->,gray] (0,-0.2)--(0,2.6);
  \draw[very thick,->,color=MidnightBlue] (0,0)--(3.4,0.55) node[right]{\small $M_1=\ker Q_1$};
  \draw[very thick,->,color=BrickRed] (0,0)--(2.0,2.3) node[above]{\small $M_2=\ker Q_2$};
  \draw (1.0,0.16) arc (9.2:49:1.0); \node at (1.25,0.62) {\small$\thetaF$};
  \node[anchor=west,align=left,font=\small] at (4.6,1.4)
   {$\cF=\cos\thetaF$, \quad $\gamma_2=1-\cF$\\[2pt]
    error amplification $\gamma_2^{-1/2}$\\[2pt]
    running example: $\thetaF\approx0.60$, $\gamma_2\approx0.17$};
\end{tikzpicture}
\caption{The Friedrichs angle $\thetaF$ between the solution spaces of two factors. A small angle
means that large pieces can almost cancel, so recovering the pieces from their sum amplifies errors by
up to $\gamma_2^{-1/2}=(1-\cos\thetaF)^{-1/2}$. The values quoted are those measured on the running
example in the companion algorithm paper.}\label{fig:friedrichs}
\end{figure}

\roleinalg{$\gN$ is computed on the sampled kernel bases and reported with every result.
It enters twice: as the error amplification $\gN^{-1/2}$ of the components, and through the Bayesian
split layer (the companion algorithm paper), whose credible intervals widen and whose split probability
returns to the prior as $\thetaF\to0$. A small angle is therefore never hidden: it shows up as wide
intervals, not as a confident but wrong split.}
